% !TeX root = main.tex
% =================================================================================
% Preamble: packages and layout. Usually, you do not need to change anything here.
% Thesis data (title, author, ...) is set in main.tex.
% =================================================================================

% ---------------------------------------------------------------------------------
% Fonts: Palatino (text and math), Helvetica (headings), Inconsolata (code).
% Works with pdfLaTeX (default) and LuaLaTeX.
% ---------------------------------------------------------------------------------
\usepackage{iftex}
\usepackage{amsmath}
\usepackage{mathtools}
\ifPDFTeX
  \usepackage[T1]{fontenc}
  \usepackage{newpxtext}
  \usepackage{newpxmath}
  \usepackage[scaled=0.90]{helvet}
  \usepackage[varqu,varl,scaled=0.95]{zi4}
\else
  \usepackage{unicode-math}
  \setmainfont{TeX Gyre Pagella}
  \setsansfont{TeX Gyre Heros}[Scale=0.90]
  \setmonofont{Inconsolatazi4}[Extension=.otf, UprightFont=*-Regular, BoldFont=*-Bold, Scale=0.95]
  \setmathfont{TeX Gyre Pagella Math}
\fi
\linespread{1.05} % Palatino needs slightly more leading
\usepackage[final]{microtype}

% ---------------------------------------------------------------------------------
% Colors, graphics, tables
% ---------------------------------------------------------------------------------
\PassOptionsToPackage{table,svgnames,dvipsnames}{xcolor}
\usepackage{xcolor}
\usepackage{graphicx}
\graphicspath{{figures/}{logos/}}
\usepackage{booktabs}    % \toprule, \midrule, \bottomrule -- avoid vertical lines in tables
\usepackage{tabularx}    % tables with a fixed total width
\usepackage{longtable}   % tables spanning multiple pages
\usepackage{multirow}    % cells spanning multiple rows
\usepackage{makecell}    % line breaks in cells: \makecell{a\\b}
\usepackage{siunitx}     % numbers and units: \num{12345.6}, \SI{3}{\milli\second}
\sisetup{group-minimum-digits=4}
\usepackage[labelfont=bf,font=small]{caption}
\usepackage{subcaption}  % sub figures: subfigure environment
\usepackage{tikz}
\usetikzlibrary{arrows.meta,positioning,shapes.geometric}
\usepackage{pgfplots}
\pgfplotsset{
  compat=1.18,
  cycle list={{TUMBlue,mark=*},{TUMAccentOrange,mark=square*},{TUMAccentGreen,mark=triangle*},{TUMSecondaryBlue2,mark=diamond*}},
}
\usepackage{pgfgantt}    % Gantt charts, e.g., for the timeline in the exposé
\usepackage{placeins}     % \FloatBarrier: place all pending floats before this point
\usepackage{pifont}
\newcommand{\cmark}{\ding{51}}
\newcommand{\xmark}{\ding{55}}

% ---------------------------------------------------------------------------------
% TUM title page, metadata, colors
% ---------------------------------------------------------------------------------
\usepackage{tumthesis}

% ---------------------------------------------------------------------------------
% Page layout and headings (KOMA-Script)
% ---------------------------------------------------------------------------------
\usepackage{scrlayer-scrpage}
\clearpairofpagestyles
\automark[section]{chapter}
\ohead{\headmark}
\ofoot*{\pagemark}
\setkomafont{pageheadfoot}{\normalfont\sffamily\small\color{TUMDarkGray}}
\setkomafont{pagenumber}{\normalfont\sffamily\small}
\pagestyle{scrheadings}

\setkomafont{disposition}{\normalcolor\sffamily\bfseries}
\addtokomafont{chapter}{\color{TUMBlue}}
\addtokomafont{chapterprefix}{\color{TUMBlue}}
\setkomafont{descriptionlabel}{\normalfont\bfseries}
\RedeclareSectionCommand[beforeskip=0pt,afterskip=1.5\baselineskip]{chapter}

% Avoid widows and orphans
\clubpenalty=10000
\widowpenalty=10000
\displaywidowpenalty=10000

% ---------------------------------------------------------------------------------
% Language-related
% ---------------------------------------------------------------------------------
\usepackage[autostyle=true]{csquotes} % \enquote{...}
\babelhyphenation[english]{
  data-base
  time-stamp
}

% ---------------------------------------------------------------------------------
% Source code: listings (no shell escape needed, works on Overleaf)
% ---------------------------------------------------------------------------------
\usepackage{scrhack} % lets listings and algorithm cooperate with KOMA-Script lists
\usepackage{listings}
\lstset{
  basicstyle=\ttfamily\small,
  keywordstyle=\bfseries\color{TUMBlue},
  commentstyle=\color{TUMGray},
  stringstyle=\color{TUMAccentOrange},
  numbers=left,
  numberstyle=\tiny\color{TUMGray},
  numbersep=8pt,
  frame=tb,
  rulecolor=\color{TUMLightGray},
  captionpos=b,
  breaklines=true,
  breakatwhitespace=true,
  showstringspaces=false,
  tabsize=2,
  columns=fullflexible,
  keepspaces=true,
  upquote=true,
  xleftmargin=2em,
  framexleftmargin=1.5em,
  literate={ä}{{\"a}}1 {ö}{{\"o}}1 {ü}{{\"u}}1 {Ä}{{\"A}}1 {Ö}{{\"O}}1 {Ü}{{\"U}}1 {ß}{{\ss}}1,
}

% ---------------------------------------------------------------------------------
% Algorithms (pseudo code)
% ---------------------------------------------------------------------------------
\usepackage[chapter]{algorithm}
\usepackage{algpseudocodex}

% Names of listings and algorithms in English and German
\makeatletter
\addto\captionsenglish{%
  \renewcommand*{\lstlistingname}{Listing}%
  \renewcommand*{\lstlistlistingname}{List of Listings}%
  \renewcommand*{\ALG@name}{Algorithm}%
  \renewcommand*{\listalgorithmname}{List of Algorithms}%
}
\addto\captionsngerman{%
  \renewcommand*{\lstlistingname}{Listing}%
  \renewcommand*{\lstlistlistingname}{Verzeichnis der Listings}%
  \renewcommand*{\ALG@name}{Algorithmus}%
  \renewcommand*{\listalgorithmname}{Verzeichnis der Algorithmen}%
}
\makeatother

% ---------------------------------------------------------------------------------
% TODO notes. For the final version, use \usepackage[disable]{todonotes}.
% ---------------------------------------------------------------------------------
\setlength{\marginparwidth}{2cm}
\usepackage[colorinlistoftodos,color=TUMAccentOrange!30,textsize=footnotesize]{todonotes}
\setuptodonotes{bordercolor=TUMAccentOrange}

% ---------------------------------------------------------------------------------
% Bibliography (biblatex + biber)
% ---------------------------------------------------------------------------------
\usepackage[
  backend=biber,
  style=numeric-comp,   % or: alphabetic, authoryear
  sorting=nyt,          % sort by name, year, title
  maxbibnames=99,       % list all authors in the bibliography
  maxcitenames=2,
  giveninits=true,
  uniquename=init,
  doi=true,
  url=true,
  isbn=false,
  eprint=false,
  urldate=long,
]{biblatex}
\addbibresource{bibliography.bib}
% Print a URL only if there is no DOI
\DeclareSourcemap{
  \maps[datatype=bibtex]{
    \map[overwrite]{
      \step[fieldsource=doi,final]
      \step[fieldset=url,null]
      \step[fieldset=urldate,null]
    }
  }
}
\usepackage{xurl} % allow line breaks in URLs everywhere

% ---------------------------------------------------------------------------------
% Hyperlinks and cross references (load late)
% ---------------------------------------------------------------------------------
\usepackage[hidelinks,unicode,bookmarksnumbered,pdfdisplaydoctitle]{hyperref}
% For colored links in the digital version, replace "hidelinks" with:
%   colorlinks, linkcolor=TUMSecondaryBlue2, citecolor=TUMBlue, urlcolor=TUMBlue
\usepackage{bookmark}
\usepackage[capitalise,nameinlink,noabbrev]{cleveref} % \cref{...} -> "Figure 2.1"
\crefname{lstlisting}{Listing}{Listings}
\Crefname{lstlisting}{Listing}{Listings}
\crefname{algorithm}{Algorithm}{Algorithms}
\Crefname{algorithm}{Algorithm}{Algorithms}

% ---------------------------------------------------------------------------------
% Acronyms: the list shows the page on which each acronym is introduced (first use).
% ---------------------------------------------------------------------------------
\usepackage[
  acronym,            % separate list of acronyms
  nomain,             % no main glossary
  toc,                % list of acronyms in table of contents
  translate=babel,
  indexonlyfirst,     % only record the page of the first use
  nopostdot,
  nogroupskip,
]{glossaries-extra}
\setabbreviationstyle[acronym]{long-short} % first use: "long form (SHORT)", afterwards: "SHORT"
\glssetcategoryattribute{acronym}{nohyperfirst}{true} % no link on first use (it is the definition)
\makenoidxglossaries % sorting by LaTeX itself -- no external tool needed

% Plain list: acronym, meaning, and dot leaders to the page of the first use (like the ToC)
\newglossarystyle{tumacronyms}{%
  \setglossarystyle{long}%
  \renewenvironment{theglossary}%
    {\setlength{\LTpre}{0pt}\setlength{\LTpost}{0pt}%
     \begin{longtable}[l]{@{}p{0.15\textwidth}@{}p{0.85\textwidth}@{}}}%
    {\end{longtable}}%
  \renewcommand*{\glossaryheader}{}%
  \renewcommand*{\glsgroupheading}[1]{}%
  \renewcommand*{\glsgroupskip}{}%
  \renewcommand*{\glossentry}[2]{%
    \glsentryitem{##1}\glstarget{##1}{\glossentryname{##1}} &
    \glossentrydesc{##1}\glspostdescription\TOCLineLeaderFill\makebox[1.6em][r]{##2}\tabularnewline}%
  \renewcommand*{\subglossentry}[3]{\glossentry{##2}{##3}}%
}
\setglossarystyle{tumacronyms}

% !TeX root = main.tex
% =================================================================================
% Acronyms -- \newacronym{<label>}{<short>}{<long>}
% Use them in the text with \gls{<label>}, \glspl{<label>} (plural), \Gls{<label>} (capitalized).
% Only acronyms that are used appear in the list, with the page of their first use.
% =================================================================================

\newacronym{ai}{AI}{artificial intelligence}
\newacronym{gqm}{GQM}{Goal Question Metric}
\newacronym{llm}{LLM}{large language model}
\newacronym{pdf}{PDF}{Portable Document Format}
\newacronym{rq}{RQ}{research question}
\newacronym{tum}{TUM}{Technical University of Munich}
\newacronym{uml}{UML}{Unified Modeling Language}


% ---------------------------------------------------------------------------------
% Lists with custom labels, e.g., \begin{enumerate}[label=(C\arabic*)]
% ---------------------------------------------------------------------------------
\usepackage{enumitem}

% ---------------------------------------------------------------------------------
% Research questions
%   \defRQ{<question>}{<label>}  states a research question (indented, numbered)
%   \refRQ{<label>}              references it: "RQ 1"
%   \discussRQ{<label>}          repeats it as a heading, e.g., in the discussion
% ---------------------------------------------------------------------------------
\newcounter{researchquestion}
\NewDocumentCommand{\defRQ}{m m}{%
  \expandafter\gdef\csname tumrqtext@#2\endcsname{#1}%
  \begin{list}{}{\setlength{\leftmargin}{2em}\setlength{\rightmargin}{2em}%
                 \setlength{\topsep}{0.5\baselineskip}}%
    \item[]\refstepcounter{researchquestion}\label{rq:#2}%
      \textbf{RQ~\theresearchquestion:}~\enquote{#1}%
  \end{list}}
\NewDocumentCommand{\refRQ}{m}{\hyperref[rq:#1]{RQ~\ref*{rq:#1}}}
\NewDocumentCommand{\discussRQ}{m}{%
  \par\addvspace{0.5\baselineskip}%
  \noindent\textbf{\refRQ{#1}:}~\emph{\csname tumrqtext@#1\endcsname}\par\nopagebreak}

% ---------------------------------------------------------------------------------
% Findings: \finding[<label>]{<text>} -> highlighted "Finding 1: <text>"
% ---------------------------------------------------------------------------------
\newcounter{finding}
\NewDocumentCommand{\finding}{o +m}{%
  \par\addvspace{0.5\baselineskip}%
  \refstepcounter{finding}\IfValueT{#1}{\label{#1}}%
  \noindent\hspace*{1.5em}{\color{TUMBlue}\vrule width 2pt}\hspace{0.8em}%
  \parbox[c]{\dimexpr\linewidth-1.5em-2pt-0.8em\relax}{%
    \strut\textbf{Finding~\thefinding:} #2\strut}%
  \par\addvspace{0.5\baselineskip}}
% ---------------------------------------------------------------------------------
% Contributions
%   \begin{contributions} \item ... \label{c:<label>} \end{contributions} -> (C1), (C2), ...
%   \refC{<label>}  references a contribution: "C1"
% ---------------------------------------------------------------------------------
\newlist{contributions}{enumerate}{1}
\setlist[contributions]{label=(C\arabic*), ref=C\arabic*, leftmargin=3em}
\NewDocumentCommand{\refC}{m}{\hyperref[c:#1]{\ref*{c:#1}}}
\crefname{contributionsi}{contribution}{contributions}
\Crefname{contributionsi}{Contribution}{Contributions}

\crefname{researchquestion}{RQ}{RQs}
\Crefname{researchquestion}{RQ}{RQs}
\crefname{finding}{Finding}{Findings}
\Crefname{finding}{Finding}{Findings}

% Abbreviations in the running text
\usepackage{xspace}
\newcommand{\eg}{e.g.,\xspace}
\newcommand{\ie}{i.e.,\xspace}
